\documentclass{article}
\usepackage[UTF8]{ctex}
\usepackage{amsmath,amssymb,amsfonts}
\usepackage{enumitem}
\usepackage{geometry}
\geometry{a4paper, margin=1in}
% === 错题本排版优化（不参与题目标识） ===
\setlength{\parindent}{0pt}
\setlength{\parskip}{0.35em plus 0.10em minus 0.08em}
\setlength{\emergencystretch}{3em}
\setlist{itemsep=0.25em, topsep=0.40em, parsep=0pt}
\allowdisplaybreaks[2]
% === 错题本排版优化结束 ===
\begin{document}
\section*{中山大学2025年高等代数真题}
\section*{143．中山大学2025年高等代数试卷}

\subsection{题目1}
1．（15 分）计算下面的行列式 \((n \geq 3)\) ：

\[
D_{n}=\left|\begin{array}{ccccc}
2 & -1 & & & \\
-1 & 2 & -1 & & \\
& \ddots & \ddots & \ddots & \\
& & -1 & 2 & -1 \\
& & & -2 & 2
\end{array}\right| .
\]

\subsection{题目2}
2．（15 分）求下列实矩阵的逆矩阵：

\[
A=\left(\begin{array}{cccc}
2 & -1 & 0 & 0 \\
-1 & 2 & -1 & 0 \\
0 & -1 & 2 & -1 \\
0 & 0 & -1 & 2
\end{array}\right)
\]

\subsection{题目3}
3．（15 分）已知线性变换 \(f: M_{2 \times 2}(\mathbb{R}) \rightarrow M_{2 \times 2}(\mathbb{R})\) 满足

\[
f(X)=F X-X F, X \in M_{2 \times 2}(\mathbb{R}) .
\]

其中 \(F=\left(\begin{array}{cc}0 & 1 \\ 0 & 0\end{array}\right)\) ．\\
（1）记 \(\alpha_{1}=\left(\begin{array}{ll}1 & 0 \\ 0 & 1\end{array}\right), \alpha_{2}=\left(\begin{array}{cc}1 & 0 \\ 0 & -1\end{array}\right), \alpha_{3}=\left(\begin{array}{ll}0 & 1 \\ 0 & 0\end{array}\right), \alpha_{4}=\left(\begin{array}{ll}0 & 0 \\ 1 & 0\end{array}\right)\) ，求 \(f\) 在 \(\alpha_{1}, \alpha_{2}, \alpha_{3}, \alpha_{4}\) 这组基下的矩阵表示。\\
（2） \(\operatorname{Ker} f\) 表示 \(f\) 的核， \(\operatorname{Im} f\) 表示 \(f\) 的值域，则 \(\operatorname{Ker} f+\operatorname{Im} f\) 是否是直和？说明理由．

\subsection{题目4}
4．（15 分）已知 \(A=\left(\begin{array}{lll}1 & -2 & 2 \\ 2 & -3 & 2 \\ 2 & -2 & 1\end{array}\right)\) ，求可逆矩阵 \(P\) ，使得 \(P^{-1} A P\) 为对角矩阵。

\subsection{题目5}
5．（15 分）记 \(\mathbb{R}[x]_{3}\) 是次数小于 3 的实系数多项式全体构成的向量空间，在 \(\mathbb{R}[x]_{3}\) 上定义双线性函数如下

\[
(f(x), g(x))=\frac{1}{2} \int_{-1}^{1} f(x) g(x) \mathrm{d} x, f(x), g(x) \in \mathbb{R}[x]_{3} .
\]

\begin{itemize}
  \item[（1）] 证明：上式定义了 \(V\) 上的内积．
  \item[（2）] 用 Schmidt 正交化方法求 \(1, x, x^{2}\) 对应的 \(V\) 的一组标准正交基。
\end{itemize}

\subsection{题目6}
6．（15 分）已知实二次型 \(Q=X^{\mathrm{T}} A X\) ，其中

\[
A=\left(\begin{array}{ccc}
4 & -2 & 2 \\
-2 & 1 & -1 \\
2 & 0 & 1
\end{array}\right)
\]

求 \(Q\) 的负惯性指数。

\subsection{题目7}
7．（15 分）已知 \(A \in M_{n}(\mathbb{C})\) 的顺序主子式均不为零，求证：存在 \(n\) 阶上三角矩阵 \(B, C\) ，使得 \(A=B^{\mathrm{T}} C\)

\subsection{题目8}
8．（15 分）已知 \(A \in M_{n}(\mathbb{C}), f(\lambda)\) 是 \(A\) 的特征多项式，\(f(\lambda)\) 与 \(f^{\prime}(\lambda)\) 的最大公因式为 \(\lambda\) ，且\(r(A)=r\left(A^{2}\right)\) ．求证：\(A\) 可相似对角化．

\subsection{题目9}
9．（15 分）已知 \(A, B \in M_{n}(\mathbb{C}), A B=B A\) ，且 \(r(A B)=r(B)\) ．求证：若 \(A\) 是幂零矩阵，则 \(B=O\) ．

\subsection{题目10}
10．（15 分）已知 \(A, B \in M_{n}(\mathbb{C})\) 均是正定矩阵，\(X \in M_{n}(\mathbb{C})\) 满足矩阵方程

\[
A X+X A=B
\]

求证：上述矩阵方程存在唯一解，并且解也是正定矩阵。
\end{document}
